% Options for packages loaded elsewhere
\PassOptionsToPackage{unicode}{hyperref}
\PassOptionsToPackage{hyphens}{url}
\PassOptionsToPackage{dvipsnames,svgnames,x11names}{xcolor}
%
\documentclass[
  11pt,
]{article}
\usepackage{amsmath,amssymb}
\usepackage{iftex}
\ifPDFTeX
  \usepackage[T1]{fontenc}
  \usepackage[utf8]{inputenc}
  \usepackage{textcomp} % provide euro and other symbols
\else % if luatex or xetex
  \usepackage{unicode-math} % this also loads fontspec
  \defaultfontfeatures{Scale=MatchLowercase}
  \defaultfontfeatures[\rmfamily]{Ligatures=TeX,Scale=1}
\fi
\usepackage{lmodern}
\ifPDFTeX\else
  % xetex/luatex font selection
\fi
% Use upquote if available, for straight quotes in verbatim environments
\IfFileExists{upquote.sty}{\usepackage{upquote}}{}
\IfFileExists{microtype.sty}{% use microtype if available
  \usepackage[]{microtype}
  \UseMicrotypeSet[protrusion]{basicmath} % disable protrusion for tt fonts
}{}
\makeatletter
\@ifundefined{KOMAClassName}{% if non-KOMA class
  \IfFileExists{parskip.sty}{%
    \usepackage{parskip}
  }{% else
    \setlength{\parindent}{0pt}
    \setlength{\parskip}{6pt plus 2pt minus 1pt}}
}{% if KOMA class
  \KOMAoptions{parskip=half}}
\makeatother
\usepackage{xcolor}
\usepackage[margin=1in]{geometry}
\usepackage{longtable,booktabs,array}
\usepackage{calc} % for calculating minipage widths
% Correct order of tables after \paragraph or \subparagraph
\usepackage{etoolbox}
\makeatletter
\patchcmd\longtable{\par}{\if@noskipsec\mbox{}\fi\par}{}{}
\makeatother
% Allow footnotes in longtable head/foot
\IfFileExists{footnotehyper.sty}{\usepackage{footnotehyper}}{\usepackage{footnote}}
\makesavenoteenv{longtable}
\setlength{\emergencystretch}{3em} % prevent overfull lines
\providecommand{\tightlist}{%
  \setlength{\itemsep}{0pt}\setlength{\parskip}{0pt}}
\setcounter{secnumdepth}{-\maxdimen} % remove section numbering
\ifLuaTeX
  \usepackage{selnolig}  % disable illegal ligatures
\fi
\IfFileExists{bookmark.sty}{\usepackage{bookmark}}{\usepackage{hyperref}}
\IfFileExists{xurl.sty}{\usepackage{xurl}}{} % add URL line breaks if available
\urlstyle{same}
\hypersetup{
  pdftitle={A proper-time dissipative term in driven superconducting qubits: a dissertation research proposal},
  pdfauthor={Trey Morris with Claude Opus 4.7},
  pdfsubject={PhD dissertation research proposal expanding the
IBM-quantum-coherence white paper (issue \#103) into a falsifiable,
multi-year program testing whether Gill's dual / proper-time
electrodynamics predicts a measurable angle- and shape-dependent
dephasing signature in driven transmon qubits.},
  colorlinks=true,
  linkcolor={blue},
  filecolor={Maroon},
  citecolor={Blue},
  urlcolor={blue},
  pdfcreator={LaTeX via pandoc}}

\title{A proper-time dissipative term in driven superconducting qubits:
a dissertation research proposal}
\author{Trey Morris with Claude Opus 4.7}
\date{2026-06-06}

\begin{document}
\maketitle

\hypertarget{a-proper-time-dissipative-term-in-driven-superconducting-qubits}{%
\section{A proper-time dissipative term in driven superconducting
qubits}\label{a-proper-time-dissipative-term-in-driven-superconducting-qubits}}

\hypertarget{abstract}{%
\subsection{Abstract}\label{abstract}}

Decoherence sets the present ceiling on superconducting quantum
computation. The dominant theoretical account of transmon dephasing
rests on the standard Maxwell field coupled to environmental noise
channels. We propose to test an alternative. Gill's dual, proper-time
reformulation of electrodynamics is mathematically equivalent to
Maxwell's theory under a change of variables, yet it carries one
structural term that the textbook theory does not\footnote{\texttt{Roadmapping/Equation\_Verification/Two\_Mathematically\_Equivalent\_Versions\_of\_Maxwells\_Equations.md}.
  Verification of Gill's dual Maxwell formulation; the dissipative third
  term with coefficient \((\mathbf{u}\!\cdot\!\mathbf{a})/b^4\) is
  recorded as Eq. (4) and recurs in the modified Liénard--Wiechert
  fields.}. The term is a first-derivative dissipative contribution
whose coefficient is the proper-time scalar
\((\mathbf{u}\!\cdot\!\mathbf{a})/b^4\). It vanishes for an
unaccelerated source, and it inherits an angular dependence through
\(\mathbf{u}\!\cdot\!\mathbf{a}=|\mathbf{u}||\mathbf{a}|\cos\theta\). A
transmon qubit driven by a shaped microwave pulse is an accelerated
source by construction. We therefore ask whether the dual term
contributes a dephasing channel with a drive-angle and pulse-shape
signature distinct from standard radiation reaction. The proposal
develops the test in three stages: a derivation of the term's
contribution to the qubit master equation, a concrete prediction for the
angle and shape dependence of \(T_2\), and an experiment on IBM Quantum
hardware. We state at the outset that the most likely outcome is a null
result, and we treat the honest specification of a falsifiable test as
the program's primary contribution.

\hypertarget{background-and-motivation}{%
\subsection{1. Background and
motivation}\label{background-and-motivation}}

The transmon is the workhorse of contemporary superconducting quantum
processors\footnote{\texttt{koch2007\_transmon}. Koch et al.,
  ``Charge-insensitive qubit design derived from the Cooper pair box,''
  Phys. Rev.~A 76, 042319 (2007).}. Its coherence time \(T_2\) limits
the depth of any circuit executed before error correction is required.
Decoherence in these devices is well characterised experimentally and is
attributed to a small set of channels: flux noise, charge noise,
dielectric loss, and quasiparticle tunnelling\footnote{\texttt{ithier2005\_decoherence}.
  Ithier et al., ``Decoherence in a superconducting quantum bit
  circuit,'' Phys. Rev.~B 72, 134519 (2005).}. Each channel is modelled
within standard electrodynamics coupled to an environmental spectral
density. The standard account is successful, and we do not dispute it.

We observe, however, that the standard account assumes the textbook
Maxwell field. Gill's proper-time formulation is an alternative to that
field theory\footnote{\texttt{Roadmapping/Equation\_Verification/Two\_Mathematically\_Equivalent\_Versions\_of\_Maxwells\_Equations.md}.
  Verification of Gill's dual Maxwell formulation; the dissipative third
  term with coefficient \((\mathbf{u}\!\cdot\!\mathbf{a})/b^4\) is
  recorded as Eq. (4) and recurs in the modified Liénard--Wiechert
  fields.}. The two formulations are mathematically equivalent under the
substitutions \(c\to b\), \(t\to\tau\), and \(\mathbf{w}\to\mathbf{u}\),
but they are not, in the sense Gill makes precise, physically
equivalent. The proper-time wave equation carries one term with no
textbook counterpart. To the extent that this term couples to the driven
condensate of a transmon, it predicts a dephasing contribution that the
standard account omits. The motivation for this proposal is that the
framework has never been tested against a condensed-matter device, and a
driven qubit is the cleanest accelerated electromagnetic source
available in the laboratory.

It is worth being precise about where a deviation could hide. The
standard channels are characterised in the idle qubit, where the source
is unaccelerated\footnote{\texttt{ithier2005\_decoherence}. Ithier et
  al., ``Decoherence in a superconducting quantum bit circuit,'' Phys.
  Rev.~B 72, 134519 (2005).}. Gate operations are different. During a
gate the drive accelerates the condensate on a timescale of tens of
nanoseconds, and it is exactly the accelerated regime in which the term
of Eq. (2.2) is nonzero. The standard noise budget is therefore measured
in the regime where the dual term vanishes by construction, and the gate
regime where it does not vanish is the regime least constrained by
existing characterisation. This is the gap the proposal targets.

\hypertarget{problem-statement-and-research-questions}{%
\subsection{2. Problem statement and research
questions}\label{problem-statement-and-research-questions}}

The dual wave equation, written for the electric field in the
proper-time formulation, takes the form

\[
\nabla^2 \mathbf{E} \;-\; \frac{1}{b^2}\,\partial_\tau^2 \mathbf{E} \;-\; \frac{\mathbf{u}\!\cdot\!\mathbf{a}}{b^4}\,\partial_\tau \mathbf{E} \;=\; \mathcal{S}, \qquad (2.1)
\]

where \(\mathcal{S}\) is the source term and \(b\) is the proper-time
speed parameter\footnote{\texttt{Roadmapping/Equation\_Verification/Two\_Mathematically\_Equivalent\_Versions\_of\_Maxwells\_Equations.md}.
  Verification of Gill's dual Maxwell formulation; the dissipative third
  term with coefficient \((\mathbf{u}\!\cdot\!\mathbf{a})/b^4\) is
  recorded as Eq. (4) and recurs in the modified Liénard--Wiechert
  fields.}. The first two terms reproduce the textbook wave equation
under the substitutions above. The third term is the structural
deviation. Its coefficient is the dissipative scalar

\[
\Gamma_{\!d} \;\equiv\; \frac{\mathbf{u}\!\cdot\!\mathbf{a}}{b^4} \;=\; \frac{|\mathbf{u}|\,|\mathbf{a}|\cos\theta}{b^4}, \qquad (2.2)
\]

in which \(\theta\) is the angle between the source velocity and its
acceleration. Equation (2.2) is the object of this proposal. It vanishes
when \(\mathbf{a}=\mathbf{0}\), and it carries an explicit angular
dependence through \(\cos\theta\).

The research questions follow directly. First, what kinematic quantities
of the driven Cooper-pair condensate play the role of \(\mathbf{u}\) and
\(\mathbf{a}\) in Eq. (2.2)? Second, does the term contribute a
dephasing rate large enough to measure at present \(T_2\) values? Third,
can the contribution's angular and shape dependence be separated
experimentally from the standard radiation-reaction term, which depends
on \(\dot{\mathbf{a}}\) rather than on \(\theta\)? These three questions
define the three phases of the program.

The symbols used throughout are collected below.

\begin{longtable}[]{@{}
  >{\raggedright\arraybackslash}p{(\columnwidth - 2\tabcolsep) * \real{0.5000}}
  >{\raggedright\arraybackslash}p{(\columnwidth - 2\tabcolsep) * \real{0.5000}}@{}}
\toprule\noalign{}
\begin{minipage}[b]{\linewidth}\raggedright
Symbol
\end{minipage} & \begin{minipage}[b]{\linewidth}\raggedright
Meaning
\end{minipage} \\
\midrule\noalign{}
\endhead
\bottomrule\noalign{}
\endlastfoot
\(b\) & proper-time speed parameter; \(b\approx c\) in vacuum. \\
\(\tau\) & proper time, the evolution parameter of the dual
formulation. \\
\(\mathbf{u},\,\mathbf{a}\) & source velocity and acceleration of the
driven condensate. \\
\(\theta\) & angle between \(\mathbf{u}\) and \(\mathbf{a}\). \\
\(\Gamma_d\) & dissipative coefficient
\((\mathbf{u}\!\cdot\!\mathbf{a})/b^4\) of Eq. (2.2). \\
\(A(\tau)\) & microwave drive envelope; \(t_g\) the gate time. \\
\(\Gamma_\varphi^{(d)}\) & proposed dual contribution to the
pure-dephasing rate. \\
\(\kappa\) & coupling constant of Eq. (3.1), fixed in Phase A. \\
\(T_2\) & transverse coherence time of the qubit. \\
\end{longtable}

\hypertarget{hypotheses-and-predictions}{%
\subsection{3. Hypotheses and
predictions}\label{hypotheses-and-predictions}}

We advance one principal hypothesis and one falsifier. The principal
hypothesis is that the dissipative term of Eq. (2.2) contributes an
additional pure-dephasing channel to the transmon master equation during
gate operations. Written as a Lindblad dephasing rate, the proposed
contribution is

\[
\Gamma_\varphi^{(d)}(\theta, A) \;=\; \kappa \int_0^{t_g} \frac{\mathbf{u}(\tau)\!\cdot\!\mathbf{a}(\tau)}{b^4}\;\big|A(\tau)\big|^2\,d\tau, \qquad (3.1)
\]

where \(A(\tau)\) is the drive envelope, \(t_g\) is the gate time, and
\(\kappa\) is a coupling constant to be fixed in Phase A. Equation (3.1)
is a proposed form, not a derived result. The derivation of \(\kappa\)
and of the correct integrand is the work of Phase A.

The falsifier is sharper. Standard Abraham--Lorentz--Dirac radiation
reaction depends on \(\dot{\mathbf{a}}\) and carries no explicit
\(\cos\theta\) dependence at fixed \(\dot{\mathbf{a}}\)\footnote{\texttt{Roadmapping/Equation\_Verification/The\_Classical\_Electron\_Problem.md}.
  Verification of Gill, Zachary \& Lindesay; proper-time treatment of
  radiation reaction, against which the dual term must be separated.}.
The dual term depends on \(\mathbf{u}\!\cdot\!\mathbf{a}\). We therefore
predict that varying the drive angle \(\theta\) at fixed
\(|\dot{\mathbf{a}}|\), fixed amplitude, and fixed gate time produces a
\(T_2\) variation of the form

\[
\frac{1}{T_2(\theta)} \;=\; \frac{1}{T_2^{(0)}} \;+\; C\,\cos\theta, \qquad (3.2)
\]

with \(C\) a constant whose sign and magnitude Phase B will predict. A
measured angular dependence of the form (3.2), not attributable to the
standard channels at fixed \(\dot{\mathbf{a}}\), supports the framework.
Its absence at the predicted sensitivity refutes the framework for this
device class.

\hypertarget{identification-of-the-condensate-kinematics}{%
\subsection{4. Identification of the condensate
kinematics}\label{identification-of-the-condensate-kinematics}}

The principal unknown is the identification of \(\mathbf{u}\) and
\(\mathbf{a}\) for a Cooper-pair condensate. We propose three candidate
identifications and will adjudicate among them in Phase A. The first is
the supercurrent velocity set by the phase gradient,

\[
\mathbf{u}_s \;=\; \frac{\hbar}{2 m_e}\,\nabla\varphi, \qquad (4.1)
\]

where \(\varphi\) is the macroscopic condensate phase\footnote{\texttt{josephson1962\_tunneling}.
  Josephson, ``Possible new effects in superconductive tunnelling,''
  Phys. Lett. 1, 251 (1962).}. The second is the AC-Josephson velocity
at the drive frequency, fixed by the Josephson relation
\(\partial_t\varphi = 2eV/\hbar\) and therefore tunable through the
drive amplitude. The third is the time derivative of the Cooper-pair
dipole moment in the transmon charge basis.

The three identifications give different effect sizes. The dimensionless
kinematic suppression carried by the term is the fourth power of the
velocity ratio,

\[
\beta^4 \;=\; \left(\frac{|\mathbf{u}|}{b}\right)^{\!4}, \qquad (4.2)
\]

evaluated at \(b\approx c\). We summarise Eq. (4.2) for each candidate,
computed explicitly rather than quoted:

\begin{longtable}[]{@{}
  >{\raggedright\arraybackslash}p{(\columnwidth - 4\tabcolsep) * \real{0.3333}}
  >{\raggedright\arraybackslash}p{(\columnwidth - 4\tabcolsep) * \real{0.3333}}
  >{\raggedright\arraybackslash}p{(\columnwidth - 4\tabcolsep) * \real{0.3333}}@{}}
\toprule\noalign{}
\begin{minipage}[b]{\linewidth}\raggedright
Identification
\end{minipage} & \begin{minipage}[b]{\linewidth}\raggedright
velocity \(|\mathbf{u}|\)
\end{minipage} & \begin{minipage}[b]{\linewidth}\raggedright
\(\beta^4\) from Eq. (4.2)
\end{minipage} \\
\midrule\noalign{}
\endhead
\bottomrule\noalign{}
\endlastfoot
Cooper-pair drift velocity & \(\sim 1~\text{m s}^{-1}\) &
\(\sim 1\times 10^{-34}\) \\
Fermi velocity (aluminium) & \(\sim 2\times 10^{6}~\text{m s}^{-1}\) &
\(\sim 2\times 10^{-9}\) \\
AC-Josephson velocity & drive-amplitude dependent & set by the drive \\
\end{longtable}

We are careful about what Eq. (4.2) does and does not say. It is the
kinematic suppression of the source term, not the suppression of the
measured dephasing rate. The map from \(\beta^4\) to a contribution to
\(T_2\) runs through the coupling \(\kappa\) of Eq. (3.1), which Phase A
must supply. No feasibility conclusion follows from \(\beta^4\) alone.
What does follow is a relative ordering: the drift-velocity
identification is suppressed by twenty-five orders of magnitude relative
to the Fermi-velocity one, so the latter and the amplitude-dependent
AC-Josephson case are the only candidates worth carrying into Phase A.

We observe that only the AC-Josephson identification offers a route in
which the drive amplitude sets the velocity, and it is the
identification on which Phase A should concentrate. The drift-velocity
candidate is excluded by Eq. (4.2) before any derivation begins.

\hypertarget{prior-work}{%
\subsection{5. Prior work}\label{prior-work}}

The transmon and its decoherence are documented in a mature literature.
Koch and co-workers introduced the transmon as a
charge-noise-insensitive design derived from the Cooper-pair
box\footnote{\texttt{koch2007\_transmon}. Koch et al.,
  ``Charge-insensitive qubit design derived from the Cooper pair box,''
  Phys. Rev.~A 76, 042319 (2007).}. Ithier and co-workers characterised
decoherence in a superconducting qubit circuit and established the
channel decomposition still in use\footnote{\texttt{ithier2005\_decoherence}.
  Ithier et al., ``Decoherence in a superconducting quantum bit
  circuit,'' Phys. Rev.~B 72, 134519 (2005).}. Krantz and co-workers
later consolidated the engineering account in a review that we take as
the standard reference for the device physics\footnote{\texttt{krantz2019\_circuit\_qed}.
  Krantz et al., ``A quantum engineer's guide to superconducting
  qubits,'' Appl. Phys. Rev.~6, 021318 (2019).}. The Josephson effect
that underlies the AC-Josephson identification of Section 4 is due to
Josephson\footnote{\texttt{josephson1962\_tunneling}. Josephson,
  ``Possible new effects in superconductive tunnelling,'' Phys. Lett. 1,
  251 (1962).}.

The character of the dominant noise is also well established. The
low-frequency dephasing of these devices is governed by noise with an
approximately \(1/f\) spectrum, whose implications for solid-state
quantum information were reviewed by Paladino and co-workers\footnote{\texttt{paladino2014\_1f\_noise}.
  Paladino, Galperin, Falci \& Altshuler, ``1/f noise: implications for
  solid-state quantum information,'' Rev.~Mod. Phys. 86, 361 (2014).}.
The coherence times set by this noise have improved steadily as
materials and geometries are refined. Place and co-workers reached
transverse times of several hundred microseconds in tantalum-based
transmons\footnote{\texttt{place2021\_tantalum\_transmon}. Place et al.,
  ``New material platform for superconducting transmon qubits with
  coherence times exceeding 0.3 milliseconds,'' Nat. Commun. 12, 1779
  (2021).}. We take these times as the sensitivity scale against which
any framework contribution must be measured. A contribution that does
not shift \(T_2\) at the few-hundred-microsecond scale is, for present
purposes, unmeasurable.

Two families of control technique are central to the experiment of
Section 6. The first is dynamical decoupling, in which timed pulse
sequences refocus low-frequency dephasing; Uhrig derived the optimal
single-qubit sequence under a sharp-cutoff bath\footnote{\texttt{uhrig2007\_dynamical\_decoupling}.
  Uhrig, ``Keeping a quantum bit alive by optimized pi-pulse
  sequences,'' Phys. Rev.~Lett. 98, 100504 (2007).}. We use such
sequences as matched baselines, because the framework's contribution is
measured as a residual against them. The second is pulse shaping, in
which the drive envelope is tailored to suppress leakage to higher
transmon levels. The derivative-removal technique of Motzoi and
co-workers is the standard tool, and it supplies the angular control of
the drive in the quadrature plane\footnote{\texttt{motzoi2009\_drag}.
  Motzoi, Gambetta, Rebentrost \& Wilhelm, ``Simple pulses for
  elimination of leakage in weakly nonlinear qubits,'' Phys. Rev.~Lett.
  103, 110501 (2009).}. The master-equation methodology rests in turn on
the theory of quantum dynamical semigroups, whose generators were
characterised by Lindblad\footnote{\texttt{lindblad1976\_generators}.
  Lindblad, ``On the generators of quantum dynamical semigroups,''
  Commun. Math. Phys. 48, 119 (1976).}. We insert the framework's
contribution as an additional Lindblad generator, which is the natural
slot for a dissipative term of the form of Eq. (3.1).

On the framework side, the dual Maxwell wave equation and its
dissipative third term are recorded in the verification of Gill's
electrodynamics paper\footnote{\texttt{Roadmapping/Equation\_Verification/Two\_Mathematically\_Equivalent\_Versions\_of\_Maxwells\_Equations.md}.
  Verification of Gill's dual Maxwell formulation; the dissipative third
  term with coefficient \((\mathbf{u}\!\cdot\!\mathbf{a})/b^4\) is
  recorded as Eq. (4) and recurs in the modified Liénard--Wiechert
  fields.}. The proper-time radiation-reaction structure, which the
falsifier of Section 3 must be separated from, is treated in the
verification of Gill's classical-electron work\footnote{\texttt{Roadmapping/Equation\_Verification/The\_Classical\_Electron\_Problem.md}.
  As {[}\^{}4{]}; the proper-time radiation-reaction structure.}. The
dual relativistic quantum mechanics that supplies the proper-time
Hamiltonian is recorded in its own verification document\footnote{\texttt{Roadmapping/Equation\_Verification/Dual\_Relativistic\_Quantum\_Mechanics\_I.md}.
  Verification of DRQM I; the proper-time Hamiltonian and dual Dirac
  apparatus.}. We note that none of these framework documents addresses
a condensed-matter device; the conceptual extension to a transmon is the
novel content of this proposal.

The methodological precedent for the test is the strong-field
radiation-reaction program. There, the experimental difficulty is not
detecting energy loss but separating a
velocity-and-acceleration-dependent reaction force from competing
effects of similar magnitude. Recent laser--electron-beam experiments
attacked exactly this separation problem\footnote{Issues \#43 and \#48,
  the proper-time radiation-reaction and bremsstrahlung prediction
  threads (github.com/temoTxt/PyPhysics/issues/43,
  github.com/temoTxt/PyPhysics/issues/48), which collect the
  strong-field radiation-reaction experiments and the framework's
  treatment of them.}. Our falsifier inherits the same logic in a
quantum-device setting: the dual term and standard radiation reaction
are of comparable structure, and the experimental content lies in
separating them through their differing dependence on the kinematic
angle. We regard the radiation-reaction experiments as the closest
existing analogue to the measurement proposed here, and we will draw on
their analysis methods for baseline subtraction.

The general question behind the proposal is how two field theories that
are mathematically equivalent, yet physically distinct, are to be told
apart in the laboratory. Gill's phrase for the relation is that the two
are mathematically equivalent but not physically equivalent. An
equivalence under a change of variables does not guarantee an
equivalence of predictions once the variables are tied to laboratory
clocks and laboratory sources. The transmon supplies a setting in which
the tie is concrete and the residual term is, in principle, switched on
by the drive.

\hypertarget{proposed-methodology}{%
\subsection{6. Proposed methodology}\label{proposed-methodology}}

The program proceeds in three phases, each with an explicit acceptance
criterion and an explicit closure condition.

\textbf{Phase A --- apparatus mapping.} We will derive the contribution
of the dissipative term to the transmon effective master equation,
starting from Eq. (2.1) and committing to one of the three
identifications of Section 4. The derivation proceeds in two steps. We
first reduce the dual field at the qubit through the proper-time
Hamiltonian of the dual relativistic quantum mechanics\footnote{\texttt{Roadmapping/Equation\_Verification/Dual\_Relativistic\_Quantum\_Mechanics\_I.md}.
  Verification of DRQM I; the proper-time Hamiltonian and dual Dirac
  apparatus.}, so that the field term of Eq. (2.2) enters as a
perturbation of the drive Hamiltonian. We then trace out the field to
obtain the induced dissipator, identifying the coefficient \(\kappa\) of
Eq. (3.1). Each algebraic step is checked symbolically in the Wolfram
language, following the campaign's verification discipline; the checks
are written as single-line evaluations so that the transport does not
silently break a multi-line expression. The numerical baseline will be
computed with \texttt{qiskit-dynamics}, whose \texttt{Solver} integrates
the Lindblad equation under a time-dependent drive envelope entered as a
\texttt{Signal}. Acceptance requires a closed-form or tractable
expression for \(\Gamma_\varphi^{(d)}(\theta, A)\) with explicit
\(b\to c\) and \(\beta\to 0\) limits, together with a numerical
\(\Delta T_2/T_2\) at IBM Heron nominal drive parameters. The closure
condition is honest and explicit. If \(\Delta T_2/T_2 < 10^{-3}\) at the
highest computable order, or if the contribution factorises into the
standard radiation-reaction account with no residual angle dependence
separable from \(\dot{\mathbf{a}}\), the program closes at Phase A.

\textbf{Phase B --- concrete predictions.} Conditional on a nonzero
Phase-A residual, we will specify drive angles and pulse-shape envelopes
that maximise and minimise the dissipative contribution at fixed drive
amplitude. Acceptance requires at least one prediction in the form of
Eq. (3.2), with a definite sign and magnitude for \(C\), compared
against state-of-the-art transmon \(T_2\) of order \(100\) to
\(300~\mu\text{s}\) and against the angular resolution of present IBM
pulse control. The predictions will be generated and stress-tested in
simulation before any hardware time is requested.

\textbf{Phase C --- IBM Quantum test.} We will pre-flight the experiment
with \texttt{qiskit-aer}, augmenting a backend noise model with a custom
Lindblad channel carrying the Phase-A operator. On-hardware execution
will use a single high-coherence transmon on Eagle- or Heron-class
hardware through \texttt{qiskit-ibm-runtime}, under the \texttt{T2Hahn}
and \texttt{T2Ramsey} protocols of \texttt{qiskit-experiments}. The
discriminating observable is the angular dependence at fixed
dynamical-decoupling sequence, measured as a residual after subtracting
matched-baseline decoupling runs. We adopt this residual construction
because CPMG and Uhrig sequences suppress low-frequency dephasing and
would otherwise mask the framework's contribution. The verdict scheme is
threefold. \textbf{Pass} --- the framework predicts the measured angle
and shape residual within the campaign's precision. \textbf{Marginal}
--- the framework matches at the precision floor only, consistent but
not discriminating. \textbf{Refuted} --- the predicted signature is
absent at the predicted sensitivity.

The angular sweep fixes the statistical design. We will sample the drive
angle at a set of values spanning \(0\) to \(\pi\), so that a
\(\cos\theta\) residual of the form (3.2) is fit against the constant
and against any even-in-\(\theta\) systematic. The detectable
coefficient \(C\) scales with the inverse square root of the total shot
count at each angle, and with the inverse of the per-point \(T_2\)
uncertainty. The Phase-B prediction for \(C\) therefore sets the
required shot budget before any hardware time is requested, and a \(C\)
below the budgeted sensitivity returns a Marginal verdict rather than a
Refutation. We will pre-register the angle set, the shot budget, and the
residual-subtraction procedure, so that the verdict is fixed by the
prediction and not by the data.

\hypertarget{preliminary-results}{%
\subsection{7. Preliminary results}\label{preliminary-results}}

The white paper that seeds this proposal establishes three results on
which the program builds\footnote{\texttt{Roadmapping/Author\_Reports/2026-05\_quantum\_coherence\_whitepaper.md}.
  The seed white paper; issue \#103,
  github.com/temoTxt/PyPhysics/issues/103.}. First, the dissipative term
of Eq. (2.2) is a verified structural feature of the dual formulation,
not a conjecture; it is recorded as Eq. (4) of the Maxwell verification
and recurs as the third term of the modified Liénard--Wiechert
fields\footnote{\texttt{Roadmapping/Equation\_Verification/Two\_Mathematically\_Equivalent\_Versions\_of\_Maxwells\_Equations.md}.
  Verification of Gill's dual Maxwell formulation; the dissipative third
  term with coefficient \((\mathbf{u}\!\cdot\!\mathbf{a})/b^4\) is
  recorded as Eq. (4) and recurs in the modified Liénard--Wiechert
  fields.}. Second, the order-of-magnitude estimates of Section 4 are in
hand, and they already exclude two of the three candidate
identifications. Third, the simulation toolchain is identified
end-to-end, so that every phase has a laptop-scale analogue before any
IBM credit is consumed.

We are careful about what these preliminary results do and do not
establish. They establish that the term exists and that a test is
specifiable. They do not establish that the term contributes measurably
to transmon dephasing; that is precisely the open question of Phase A.

The precision-atomic campaign supplies the reason to look at a
condensed-matter device at all. In the bound-electron g-factor and the
hydrogenic spectra, the framework reproduces the textbook
quantum-electrodynamic result by construction; the cutoff carries no
nuclear-charge dependence, and the framework inherits the bound-state
structure of standard theory rather than predicting a deviation from
it\footnote{Issues \#78 and \#82, the Li\(^{2+}\) Z-extension and
  hydrogenic g-factor Z-scan (github.com/temoTxt/PyPhysics/issues/78,
  github.com/temoTxt/PyPhysics/issues/82); the Z-scan returned Outcome
  C, in which the cutoff carries no Z-dependence and the framework
  inherits the bound-state quantum-electrodynamic structure.}. A test
that the framework passes by construction cannot discriminate. The
driven condensate is attractive precisely because the dissipative term
is switched on by acceleration, a regime the atomic campaign never
probes. The proposal is therefore the campaign's first attempt at a
setting where the framework could, in principle, depart from standard
theory rather than reproduce it.

\hypertarget{timeline-and-milestones}{%
\subsection{8. Timeline and milestones}\label{timeline-and-milestones}}

We propose a four-year program. Year one is devoted to Phase A: the
symbolic derivation, the Wolfram verification, and the
\texttt{qiskit-dynamics} baseline, ending in the Phase-A acceptance or
closure decision. Year two, conditional on a nonzero residual, develops
the Phase-B predictions and the full \texttt{qiskit-aer} simulation
sweep over drive angle and pulse shape. Year three executes Phase C on
IBM Quantum hardware and performs the residual analysis against
matched-decoupling baselines. Year four is reserved for analysis,
replication, and dissertation writing. We note that the Phase-A closure
condition may terminate the program at the end of year one; in that case
the dissertation documents the closure as the framework's terminal
verdict for this device class, and the remaining years redirect to a
fallback question identified in Section 10.

The milestones and their deliverables are as follows.

\begin{longtable}[]{@{}
  >{\raggedright\arraybackslash}p{(\columnwidth - 6\tabcolsep) * \real{0.2500}}
  >{\raggedright\arraybackslash}p{(\columnwidth - 6\tabcolsep) * \real{0.2500}}
  >{\raggedright\arraybackslash}p{(\columnwidth - 6\tabcolsep) * \real{0.2500}}
  >{\raggedright\arraybackslash}p{(\columnwidth - 6\tabcolsep) * \real{0.2500}}@{}}
\toprule\noalign{}
\begin{minipage}[b]{\linewidth}\raggedright
Year
\end{minipage} & \begin{minipage}[b]{\linewidth}\raggedright
Phase
\end{minipage} & \begin{minipage}[b]{\linewidth}\raggedright
Milestone deliverable
\end{minipage} & \begin{minipage}[b]{\linewidth}\raggedright
Gate
\end{minipage} \\
\midrule\noalign{}
\endhead
\bottomrule\noalign{}
\endlastfoot
1 & A & Wolfram-verified expression for
\(\Gamma_\varphi^{(d)}(\theta,A)\); \texttt{qiskit-dynamics} baseline
\(\Delta T_2/T_2\). & Accept if residual \(\geq 10^{-3}\); else
close. \\
2 & B & Pre-registered \(\cos\theta\) prediction (3.2) with signed
\(C\); full \texttt{qiskit-aer} angle/shape sweep. & Accept if \(C\)
exceeds budgeted sensitivity. \\
3 & C & Hardware angle sweep on Eagle/Heron; residual after matched
decoupling. & Pass / Marginal / Refuted per Section 6. \\
4 & --- & Replication, analysis, dissertation. & Defence. \\
\end{longtable}

The gate column makes each year's continuation conditional on the prior
year's result, so that a null finding closes the program cleanly rather
than leaving it open-ended.

\hypertarget{expected-contributions-and-broader-impact}{%
\subsection{9. Expected contributions and broader
impact}\label{expected-contributions-and-broader-impact}}

The program's primary contribution is the specification of a falsifiable
test that does not presently exist in the framework's repertoire. At
present no condensed-matter prediction of the dual framework is on the
books. We regard the test as worth specifying whether or not the
framework passes it. A secondary contribution, independent of the
framework's fate, is a reusable \texttt{qiskit-dynamics} and
\texttt{qiskit-experiments} pipeline for inserting a custom Lindblad
channel and measuring its angular signature as a decoupling residual;
this pipeline is of use to any program that posits a non-standard
dephasing channel. The broader impact is methodological. A clean
experimental adjudication of a framework that is mathematically
equivalent to Maxwell's theory, yet physically distinct, would sharpen
the general question of how such equivalences are to be tested.

\hypertarget{risks-feasibility-and-alternatives}{%
\subsection{10. Risks, feasibility, and
alternatives}\label{risks-feasibility-and-alternatives}}

The principal risk is that the effect is unmeasurably small. The
order-of-magnitude table of Section 4 makes this risk explicit, and the
Phase-A closure condition converts it into a clean decision rather than
an open-ended search. A second risk is that the dual contribution is not
separable from standard radiation reaction; the falsifier of Eq. (3.2)
is constructed precisely to test separability, and a negative
separability result is itself a publishable closure. A third risk is
access: pulse-level control on IBM hardware may not be provisioned, in
which case gate-level DRAG-coefficient sweeps provide a coarser but
adequate angular handle.

The fallback question, should Phase A close the primary program, is
whether the same dissipative term contributes to the coherence of
trapped-ion or neutral-atom qubits, where the accelerated source is a
single charged particle rather than a condensate and the identification
of \(\mathbf{u}\) and \(\mathbf{a}\) is unambiguous. This fallback
inherits the entire methodology of Sections 6 through 8 and requires
only a change of platform. We regard the fallback as a strength of the
proposal, because the central derivation is platform-independent.

\hypertarget{resources-required}{%
\subsection{11. Resources required}\label{resources-required}}

The program is modest in physical resources and concentrated in
analytical effort. Years one and two require only commodity computing
for the Wolfram verification and the \texttt{qiskit-dynamics} and
\texttt{qiskit-aer} simulations, both of which run at laptop scale. Year
three requires IBM Quantum Network access to a single Eagle- or
Heron-class device, with pulse-level control through
\texttt{qiskit.pulse} if it can be provisioned and gate-level
DRAG-coefficient sweeps as the fallback. The shot budget is fixed by the
Phase-B prediction for the coefficient \(C\), and credit is requested
only after that prediction is in hand. The advisory requirement is the
one genuine dependency: the apparatus-mapping of Section 4 needs a
superconducting-qubit theorist to confirm the condensate kinematics, and
the framework-internal derivation of Phase A benefits from contact with
the framework's authors. We treat both as collaborations to be secured
in year one, not as open risks to be discovered later.

\hypertarget{appendix-a-phase-a-derivation-scaffold-for-review}{%
\subsection{Appendix A --- Phase-A derivation scaffold (for
review)}\label{appendix-a-phase-a-derivation-scaffold-for-review}}

This appendix sets out the algebra of Phase A as far as standard
open-quantum-systems theory carries it, and it marks the points at which
the framework supplies an input that the standard theory does not. It is
a scaffold, not a result. The standard steps, A.1 and A.4, are written
out. The framework-specific steps, A.2, A.3, and A.5, carry the
substantive assumptions; the coupling \(\kappa\) is left symbolic
throughout, because fixing it is precisely the work of Phase A. We write
the appendix in this form so that the framework's authors can check the
framework-specific steps directly, independently of the standard
machinery around them.

\hypertarget{a.1-the-driven-transmon-standard}{%
\subsubsection{A.1 The driven transmon
(standard)}\label{a.1-the-driven-transmon-standard}}

We treat the transmon as a weakly anharmonic oscillator, in the manner
of Koch and co-workers\footnote{\texttt{koch2007\_transmon}. Koch et
  al., ``Charge-insensitive qubit design derived from the Cooper pair
  box,'' Phys. Rev.~A 76, 042319 (2007).}. The bare Hamiltonian is

\[
H_0/\hbar \;=\; \omega_q\,a^\dagger a \;+\; \tfrac{\alpha}{2}\,a^\dagger a^\dagger a\,a, \qquad (A.1)
\]

with \(\omega_q\) the qubit frequency and \(\alpha\) the anharmonicity.
A shaped microwave drive couples through the charge operator,

\[
H_{\rm drive}(t)/\hbar \;=\; \Omega(t)\,\big[\cos\phi\,(a+a^\dagger) + \sin\phi\;i(a^\dagger-a)\big]\cos\omega_d t, \qquad (A.2)
\]

where \(\Omega(t)\propto A(\tau)\) is the envelope and \(\phi\) is the
phase in the in-phase/quadrature plane. The phase \(\phi\) is the
experimentally controlled quantity. It is the handle through which the
drive angle of Section 3 is varied, and the derivative-removal shaping
of Motzoi and co-workers operates on \(\Omega(t)\) at fixed
\(\phi\)\footnote{\texttt{motzoi2009\_drag}. Motzoi, Gambetta,
  Rebentrost \& Wilhelm, ``Simple pulses for elimination of leakage in
  weakly nonlinear qubits,'' Phys. Rev.~Lett. 103, 110501 (2009).}.

\hypertarget{a.2-the-dual-term-as-a-friction-on-the-field-mode-framework}{%
\subsubsection{A.2 The dual term as a friction on the field mode
(framework)}\label{a.2-the-dual-term-as-a-friction-on-the-field-mode-framework}}

We expand the field in modes
\(\mathbf{E}(\mathbf{x},\tau)=\sum_{\mathbf k}E_{\mathbf k}(\tau)\,e^{i\mathbf{k}\cdot\mathbf{x}}\)
and substitute into the dual wave equation (2.1). Each mode then obeys

\[
\partial_\tau^2 E_{\mathbf k} \;+\; \Gamma_{\rm mode}\,\partial_\tau E_{\mathbf k} \;+\; (b\,|\mathbf{k}|)^2 E_{\mathbf k} \;=\; -\,b^2\,\mathcal{S}_{\mathbf k}, \qquad \Gamma_{\rm mode}=\frac{\mathbf{u}\!\cdot\!\mathbf{a}}{b^2}. \qquad (A.3)
\]

We observe that the dual term enters Eq. (A.3) as a damping of the field
mode. The modal damping rate \(\Gamma_{\rm mode}\) has dimensions of
inverse time, and it is nonzero only while the source accelerates. This
is the structural content of the framework at this order: the
proper-time formulation makes the electromagnetic mode lossy in
proportion to \(\mathbf{u}\!\cdot\!\mathbf{a}\), where the textbook
formulation makes it lossless.

\hypertarget{a.3-the-condensate-kinematics-under-drive-framework-gated}{%
\subsubsection{A.3 The condensate kinematics under drive (framework,
gated)}\label{a.3-the-condensate-kinematics-under-drive-framework-gated}}

We adopt the AC-Josephson identification of Section 4, in which the
junction phase obeys the Josephson relation
\(\partial_\tau\varphi = 2eV(\tau)/\hbar\) with
\(V(\tau)\propto A(\tau)\). The condensate velocity is then set by the
drive, and its acceleration is the proper-time derivative of that
velocity. Writing \(u(\tau)=|\mathbf{u}(\tau)|\) and
\(a(\tau)=|\mathbf{a}(\tau)|\), Eq. (A.3) becomes

\[
\Gamma_{\rm mode}(\tau) \;=\; \frac{u(\tau)\,a(\tau)\,\cos\theta}{b^2}. \qquad (A.4)
\]

The angle \(\theta\) enters through \(\mathbf{u}\!\cdot\!\mathbf{a}\).
Its relation to the experimentally controlled phase \(\phi\) of Eq.
(A.2) is the mapping that Phase A must establish; we do not assume
\(\theta=\phi\).

\hypertarget{a.4-systembath-reduction-standard}{%
\subsubsection{A.4 System--bath reduction
(standard)}\label{a.4-systembath-reduction-standard}}

We write the qubit coupled to the electromagnetic modes as
\(H = H_q + H_B + H_{\rm int}\), with a linear coupling
\(H_{\rm int}=\hbar\,(a+a^\dagger)\otimes B\) to a bath operator \(B\)
built from the mode amplitudes \(E_{\mathbf k}\). The standard
Born--Markov--secular reduction then yields a master equation of
Lindblad form\footnote{\texttt{lindblad1976\_generators}. Lindblad, ``On
  the generators of quantum dynamical semigroups,'' Commun. Math. Phys.
  48, 119 (1976).},

\[
\dot\rho \;=\; -\tfrac{i}{\hbar}[H_q,\rho] \;+\; \gamma_\downarrow\,\mathcal{D}[a]\rho \;+\; \gamma_\uparrow\,\mathcal{D}[a^\dagger]\rho \;+\; \gamma_\varphi\,\mathcal{D}[a^\dagger a]\rho, \qquad (A.5)
\]

with the dissipator
\(\mathcal{D}[L]\rho = L\rho L^\dagger - \tfrac12\{L^\dagger L,\rho\}\).
Pure dephasing is the \(\mathcal{D}[a^\dagger a]\) channel, and its rate
is fixed by the bath spectral density at zero frequency,

\[
\gamma_\varphi \;=\; \tfrac12\,S_B(\omega\to 0), \qquad S_B(\omega)=\int_{-\infty}^{\infty}\!\langle B(\tau)B(0)\rangle\,e^{i\omega\tau}\,d\tau. \qquad (A.6)
\]

Equation (A.6) is the slot into which the framework's contribution must
be inserted. A drive-induced, angle-dependent modification of the bath
correlation \(\langle B(\tau)B(0)\rangle\) produces a corresponding
modification of \(\gamma_\varphi\).

\hypertarget{a.5-the-frameworks-contribution-to-the-dephasing-rate-the-key-step-gated}{%
\subsubsection{A.5 The framework's contribution to the dephasing rate
(the key step,
gated)}\label{a.5-the-frameworks-contribution-to-the-dephasing-rate-the-key-step-gated}}

The modal friction of Eq. (A.3) modifies the field correlations that the
qubit samples. The bath operator \(B\) inherits the loss
\(\Gamma_{\rm mode}\), so the dephasing spectral density acquires an
extra piece. Carrying that piece through Eq. (A.6) and averaging over
the gate gives the proposed contribution,

\[
\gamma_\varphi^{(d)}(\theta,A) \;=\; \kappa \int_0^{t_g} \frac{u(\tau)\,a(\tau)\,\cos\theta}{b^4}\,\big|A(\tau)\big|^2\,d\tau, \qquad (A.7)
\]

which is Eq. (3.1) of the body. The constant \(\kappa\) collects the
qubit--mode coupling, the mode density of states at \(\omega_q\), and
the powers of \(b\) that convert the field-equation coefficient of Eq.
(2.2) into a rate. We note one bookkeeping point for the reviewer. The
modal rate of Eq. (A.3) carries \(1/b^2\), while Eq. (A.7) is written
with \(1/b^4\) to match the body; the difference is absorbed into
\(\kappa\), and the explicit powers of \(b\) must be tracked through the
reduction rather than assumed. The value of \(\kappa\) is the principal
output of Phase A and is not fixed by this scaffold.

\hypertarget{a.6-limits-and-the-falsifier}{%
\subsubsection{A.6 Limits and the
falsifier}\label{a.6-limits-and-the-falsifier}}

Two checks constrain any completed derivation. The slow-source limit is

\[
\beta\to 0 \;\Rightarrow\; u\to 0 \;\Rightarrow\; \Gamma_{\rm mode}\to 0 \;\Rightarrow\; \gamma_\varphi^{(d)}\to 0, \qquad (A.8)
\]

so the framework contribution vanishes and the standard account is
recovered. The textbook limit \(b\to c\) leaves Eq. (A.7) finite and is
the regime of the experiment. The falsifier of Section 3 follows from
the angular structure of Eq. (A.7). Standard radiation reaction
contributes a rate proportional to \(\dot{\mathbf a}\) with no explicit
\(\cos\theta\) at fixed \(\dot{\mathbf a}\). Holding
\(|\dot{\mathbf a}|\), the amplitude, and the gate time fixed while
varying \(\theta\) therefore isolates the framework term,

\[
\frac{1}{T_2(\theta)} \;=\; \frac{1}{T_2^{(0)}} \;+\; C\,\cos\theta, \qquad C=\frac{\kappa}{b^4}\Big\langle u\,a\,|A|^2\Big\rangle t_g, \qquad (A.9)
\]

which is Eq. (3.2) with \(C\) now expressed through the same \(\kappa\)
as Eq. (A.7). A measured \(\cos\theta\) residual fixes the magnitude and
sign of \(C\), and through it tests the framework's prediction for
\(\kappa\).

\hypertarget{a.7-open-points-for-review}{%
\subsubsection{A.7 Open points for
review}\label{a.7-open-points-for-review}}

The scaffold leaves five points open. We list them so that the review
can target them directly.

\begin{enumerate}
\def\labelenumi{\arabic{enumi}.}
\tightlist
\item
  The identification of \(\mathbf{u}\) and \(\mathbf{a}\) for the
  condensate (A.3), and in particular whether the AC-Josephson velocity
  is the correct choice in the transmon charge basis.
\item
  The map from the experimentally controlled IQ phase \(\phi\) of Eq.
  (A.2) to the geometric angle \(\theta\) of Eq. (A.4).
\item
  The propagation of the modal friction \(\Gamma_{\rm mode}\) into the
  dephasing spectral density of Eq. (A.6), which is the load-bearing
  framework step.
\item
  The value and sign of \(\kappa\), and the explicit powers of \(b\)
  that the reduction carries.
\item
  The reconciliation of the kinematic suppression \(\beta^4\) of Eq.
  (4.2) with the rate-level bookkeeping of Eqs. (A.3) and (A.7).
\end{enumerate}

We regard points 3 and 4 as the ones requiring the framework's authors.
Points 1, 2, and 5 can be settled with a superconducting-qubit theorist
and the campaign's verification tooling.

\hypertarget{references}{%
\subsection{References}\label{references}}

\end{document}
